\begin{center}
  \large\textbf{ABSTRAK}
\end{center}

\addcontentsline{toc}{chapter}{ABSTRAK}

\vspace{2ex}

\begingroup
% Menghilangkan padding
\setlength{\tabcolsep}{0pt}

\noindent
\begin{tabularx}{\textwidth}{l >{\centering}m{2em} X}
  Nama Mahasiswa    & : & \name{}         \\

  Judul Tugas Akhir & : & \tatitle{}      \\

  Pembimbing        & : & 1. \advisor{}   \\
                    &   & 2. \coadvisor{} \\
\end{tabularx} 
\endgroup

% Ubah paragraf berikut dengan abstrak dari tugas akhir
Penelitian ini merancang dan mengevaluasi sistem \textit{Credit Risk Scoring} berbasis \textit{Machine Learning Operations} (MLOps) untuk mendukung penilaian risiko kredit yang adaptif, terukur, dan siap dioperasikan secara berkelanjutan. Dataset sintetis sebanyak 10.000 sampel digunakan untuk merepresentasikan profil nasabah kredit ritel, kemudian diproses melalui transformasi \textit{Weight of Evidence} (WoE), seleksi fitur berbasis \textit{Information Value} (IV), dan penanganan ketidakseimbangan kelas. Model utama dibangun menggunakan XGBoost, sedangkan pipeline MLOps mengintegrasikan pelacakan eksperimen, registrasi model, deployment API, serta otomasi pelatihan ulang (\textit{Continuous Training}) yang dieksekusi secara terjadwal maupun berbasis \textit{event} (akumulasi data baru). Hasil pengujian menunjukkan bahwa pendekatan ini mampu menghasilkan model yang akurat, lebih mudah diinterpretasikan melalui SHAP, serta didukung alur otomasi yang secara konsisten memastikan kebaruan model (\textit{model freshness}) guna meningkatkan keandalan implementasi pada lingkungan lembaga keuangan.

% Ubah kata-kata berikut dengan kata kunci dari tugas akhir
Kata Kunci: credit risk scoring, MLOps, XGBoost, WoE, SHAP
